% ============================================================
% Notes: Nonlocality + Volterra Principles + Nonlocal Operators
%        + Relation to Fractional Calculus
% Source: Richard Herrmann, Fractional Calculus (World Scientific)
%         esp. Ch. 1 (Intro) and Ch. 8 (Nonlocality and Memory Effects)
% ============================================================

\section{Concept of nonlocality}

\subsection{Local vs nonlocal description}
A \textbf{local} theory determines the evolution at $(x,t)$ using only values at that point
(and finitely many derivatives):
\begin{equation}
\mathcal{E}\big(u(x,t),\partial u(x,t),\partial^2 u(x,t),\dots\big)=0.
\end{equation}

A \textbf{nonlocal} theory allows the state at $(x,t)$ to depend on values in a region
of space and/or the history in time:
\begin{equation}
u(x,t) = \int_{\Omega}K(x,\xi)\,g(\xi,t)\,d\xi
\qquad \text{(spatial nonlocality)},
\end{equation}
\begin{equation}
u(x,t) = \int_0^t M(t-\tau)\,g(x,\tau)\,d\tau
\qquad \text{(temporal nonlocality / memory)}.
\end{equation}

\paragraph{Key physical interpretation (Herrmann).}
Nonlocal effects may occur in \textbf{space and time}. In the time domain, nonlocality
appears as a \textbf{memory effect}: the current behavior depends not only on the present
state but also on the past. :contentReference[oaicite:1]{index=1}

% ------------------------------------------------------------

\section{Volterra principles of nonlocality}

\subsection{Volterra (hereditary) principle}
Volterra's viewpoint: in systems with memory, the response at time $t$ depends on
the full past history through a \textbf{causal integral operator}:
\begin{equation}
R(t)=R_0(t)+\int_0^t K(t,\tau)\,I(\tau)\,d\tau,
\qquad (t\ge 0).
\end{equation}
Here $I(\tau)$ is the ``input'' (excitation/forcing), $R(t)$ the response, and $K$ is a memory kernel.

\subsection{Volterra superposition principle (linear memory)}
If the system is linear, contributions from different past times add:
\begin{equation}
R(t) = \int_0^t M(t-\tau)\,I(\tau)\,d\tau,
\end{equation}
i.e. a \textbf{Volterra convolution} with a kernel $M$ depending only on time difference.

\subsection{Causality constraint}
The Volterra form enforces causality automatically:
\begin{equation}
\tau>t \ \Rightarrow\  \text{no influence on }R(t),
\qquad \text{hence integration domain }[0,t].
\end{equation}

\subsection{Geometric memory effect (Herrmann example)}
Herrmann gives a simple geometric mechanism: adding boundaries to an otherwise free system
can produce a delayed contribution to the source term, e.g.
\begin{equation}
\psi(t)\quad \longrightarrow\quad \psi(t)+\rho\,\psi(t-\tau),
\end{equation}
where the second term is a delayed (nonlocal-in-time) contribution generated by reflections. :contentReference[oaicite:2]{index=2}

\paragraph{Important point.}
The system does not ``store intelligence''; memory is an emergent effect from the dynamics
(e.g. reflections/boundaries) rather than literal memorization. :contentReference[oaicite:3]{index=3}

% ------------------------------------------------------------

\section{Nonlocal operations (nonlocal operators)}

\subsection{Spatial nonlocal operator (integral kernel)}
A general spatial nonlocal operator is
\begin{equation}
(\mathcal{K}u)(x)=\int_{\Omega}K(x,\xi)\,u(\xi)\,d\xi.
\end{equation}
This contrasts with local operators like $\partial_x u(x)$ which only use infinitesimal neighborhood information.

\subsection{Translation-invariant kernels $\Rightarrow$ convolution}
If $K(x,\xi)=K(x-\xi)$, then
\begin{equation}
(\mathcal{K}u)(x)=(K*u)(x)=\int_{\mathbb{R}}K(x-\xi)\,u(\xi)\,d\xi.
\end{equation}

\subsection{Temporal nonlocal operator (memory kernel)}
A general causal memory operator is
\begin{equation}
(\mathcal{M}u)(t)=\int_0^t M(t-\tau)\,u(\tau)\,d\tau.
\end{equation}

\subsection{Kernel types (how memory/range behaves)}
\begin{itemize}
\item \textbf{Short memory:} $M(t)\sim e^{-t/\tau}$ (finite relaxation scale).
\item \textbf{Long memory:} $M(t)\sim t^{-\alpha}$ (power-law tail; history never fully dies).
\item \textbf{Local limit:} $M(t)\to \delta(t)$ gives $(\mathcal{M}u)(t)=u(t)$.
\end{itemize}

% ------------------------------------------------------------

\section{Relation to fractional calculus}

\subsection{Fractional calculus as a subclass of nonlocal theories}
Herrmann emphasizes that an important subgroup of nonlocal field theories can be described
by \textbf{fractional operators}, i.e. derivatives of non-integer order:
\begin{equation}
\frac{d^n}{dx^n}\ \longrightarrow\ \frac{d^\alpha}{dx^\alpha},
\qquad \alpha\in\mathbb{R},\mathbb{C}\ (\text{or even }\alpha=\alpha(x,t)). :contentReference[oaicite:4]{index=4}
\end{equation}

\subsection{Why fractional operators are nonlocal}
Fractional derivatives typically involve integrals with weakly singular kernels.
This means the derivative at $t$ depends on values for all $\tau<t$ (memory).

\subsection{Fractional integral (Volterra operator)}
For $\alpha>0$, the Riemann--Liouville fractional integral is
\begin{equation}
(I^\alpha f)(t)
=
\frac{1}{\Gamma(\alpha)}
\int_0^t (t-\tau)^{\alpha-1}f(\tau)\,d\tau.
\end{equation}
This is exactly a Volterra convolution:
\[
(I^\alpha f)(t) = (M_\alpha * f)(t),
\qquad
M_\alpha(t)=\frac{t^{\alpha-1}}{\Gamma(\alpha)}.
\]

\subsection{Fractional derivatives (Caputo vs Riemann--Liouville)}
For $0<\alpha<1$:

\paragraph{Riemann--Liouville derivative.}
\begin{equation}
(D_t^\alpha f)(t)
=
\frac{d}{dt}
\left[
\frac{1}{\Gamma(1-\alpha)}
\int_0^t (t-\tau)^{-\alpha}f(\tau)\,d\tau
\right].
\end{equation}

\paragraph{Caputo derivative.}
\begin{equation}
({}^{C}D_t^\alpha f)(t)
=
\frac{1}{\Gamma(1-\alpha)}
\int_0^t (t-\tau)^{-\alpha}f'(\tau)\,d\tau.
\end{equation}

\paragraph{Difference (important in physics).}
Caputo gives ${}^{C}D^\alpha(\text{const})=0$ while Riemann--Liouville does not.
Herrmann explicitly highlights this as a motivation for Caputo-type definitions. :contentReference[oaicite:5]{index=5}

\subsection{Fractional derivative defined via eigenfunctions (Herrmann's ``historical'' route)}
Herrmann presents an intuitive definition based on preserving the integer-order derivative rules
for special function classes (then extending $n\to \alpha$). For example: :contentReference[oaicite:6]{index=6}
\begin{equation}
\frac{d^\alpha}{dx^\alpha}e^{kx}=k^\alpha e^{kx},
\qquad k\ge 0,
\end{equation}
\begin{equation}
\frac{d^\alpha}{dx^\alpha}\sin(kx)
=
k^\alpha\sin\!\left(kx+\frac{\pi\alpha}{2}\right),
\qquad k\ge 0,
\end{equation}
\begin{equation}
\frac{d^\alpha}{dx^\alpha}x^k
=
\frac{\Gamma(1+k)}{\Gamma(1+k-\alpha)}x^{k-\alpha},
\qquad x\ge 0.
\end{equation}

\subsection{Nonlocality via infinite-derivative expansion}
Herrmann also shows that fractional differentiation can be written as an infinite series
in integer derivatives (binomial-type expansion):
\begin{equation}
\left(\frac{d}{dx}+\omega\right)^\alpha
=
\sum_{k=0}^{\infty}\binom{\alpha}{k}\omega^{\alpha-k}\frac{d^k}{dx^k}.
\end{equation}
Since infinitely many derivatives are required, one needs essentially the full Taylor data of $f$,
which is another way to see \textbf{nonlocality} (global information requirement). :contentReference[oaicite:7]{index=7}

\subsection{Memory kernels $\Rightarrow$ fractional dynamics (core bridge)}
Volterra memory:
\[
R(t)=\int_0^t M(t-\tau)I(\tau)\,d\tau
\]
becomes fractional when the kernel has a power-law form:
\begin{equation}
M(t)\propto t^{-\alpha}
\quad\Rightarrow\quad
R(t)\propto {}^{C}D_t^\alpha I(t)\ \text{or}\ I^\alpha I(t),
\end{equation}
depending on the constitutive structure.

\paragraph{Interpretation.}
\begin{itemize}
\item Exponential kernels $\Rightarrow$ short memory, Markov-like relaxation.
\item Power-law kernels $\Rightarrow$ long memory, fractional operators.
\end{itemize}

% ------------------------------------------------------------

\section{High-yield summary}

\begin{itemize}
\item \textbf{Nonlocality} = dependence on a region in space or the entire past in time.
\item \textbf{Volterra nonlocality} = causal memory integral on $[0,t]$:
\[
R(t)=\int_0^t M(t-\tau)I(\tau)\,d\tau.
\]
\item \textbf{Nonlocal operators} are integral-kernel operators in space/time.
\item \textbf{Fractional calculus} is a structured way to build nonlocal operators:
\[
\frac{d^n}{dx^n}\to \frac{d^\alpha}{dx^\alpha}.
\]
\item \textbf{Fractional integrals/derivatives} are nonlocal because they use Volterra-type
power-law kernels:
\[
(I^\alpha f)(t)=\frac{1}{\Gamma(\alpha)}\int_0^t (t-\tau)^{\alpha-1}f(\tau)\,d\tau.
\]
\item Herrmann's key message: fractional operators provide a practical framework for
nonlocal theories with memory effects. :contentReference[oaicite:8]{index=8}
\end{itemize}
